\section{Counting equivalence at every finite rank}
\label{hf:section}
\label{sec:hf-transfer}

Small supports connect the atom structure with the sets built during a
choiceless computation.  The connection must preserve exact counts, even
though supported objects can have arbitrarily large finite rank.
We prove an abstract transfer theorem for this purpose.  Its proof first
represents supported objects by finite expressions over support tuples,
then compares the finite orbits of those objects after parameters have
been fixed.  Matching these orbits will produce a bijection of the entire
supported domains.
Forms and molecules originate in Blass, Gurevich and Shelah
\cite[Sections~8--9, published version]{bgs1999} and were developed
for counting by Dawar, Richerby and Rossman
\cite[Section~8, Theorem~33, author version]{drr2008}.
Their hereditary-support transfer already permits arbitrary finite rank.
We give the quantitative group-relative version proved in
\cite[Section~5, Theorem~14]{cptCompanion2026}, including its complete
orbit-counting argument.

Let $B$ be a finite relational structure, also writing $B$ for its
universe, and let $K\leq\Aut(B)$.  Regard the elements of $B$ as atoms,
distinct from every set, and extend the action of $K$ recursively to
$\HF(B)$.  A tuple $\alpha$ of atoms \emph{supports} $x\in\HF(B)$ if
every element of $K$ fixing $\alpha$ coordinatewise fixes $x$.
For an integer $s\geq1$, let
\[
 \begin{aligned}
 D_s(B,K)=\{x\in\HF(B):{}&
 \text{$x$ and every recursive member of $x$}\\
 &\text{have a $K$-support of length at most $s$}\}.
 \end{aligned}
\]
Here a recursive member is an object reached by a positive number of
membership steps.  We view $D_s(B,K)$ as a relational structure with
membership, a unary atom predicate, and the relations of $B$, interpreted
as false whenever an argument is a set.  This domain is closed under
taking members and is invariant under $K$.  It contains every atom and
every pure hereditarily finite set, since an atom supports itself and a
pure set has empty support.  Consequently the domain is countably
infinite when $B$ is finite, although each of its individual objects is
hereditarily finite.

For a tuple $\bar a$ of length at most $M$, write
$\operatorname{tp}^{B}_M(\bar a)$ for its complete $C^M$ type:
the set of all $C^M$ formulas true at $\bar a$, with free variables among
the indicated ordered tuple positions.  Tuple entries may repeat.
The variable bound restricts the stock of variable names, not the length
of a formula or the natural-number constants in exact counting
quantifiers.

\begin{theorem}[Transfer through small supports]
\label{hf:transfer}
Let $B_0,B_1$ be nonempty finite structures in the same finite relational
vocabulary of arity at most two, and let
$K_i\leq\Aut(B_i)$ for $i=0,1$.
Let $s,m,M$ be positive integers with
\[
                         \max\{4,m+1\}s\leq M.
\]
Suppose that:
\begin{enumerate}
\item $B_0$ and $B_1$ satisfy the same $C^M$ sentences;
\item in each $B_i$, any two tuples of length at most
      $\max\{2,m\}s$ having the same $C^M$ type belong to the same
      $K_i$-orbit.
\end{enumerate}
Then $D_s(B_0,K_0)$ and $D_s(B_1,K_1)$ satisfy the same $C^m$ sentences
in membership, atomhood, and the base relations.
\end{theorem}

All supports below are padded to length exactly $s$.
This is possible by repeating entries, or by using an arbitrary atom
when the support is empty; the nonemptiness hypothesis is used here.
Padding preserves the support property.  Write
$D_i=D_s(B_i,K_i)$ for the rest of the proof.

\subsection{Exact extension counts of atom types}

Only types realized in $B_0$ or $B_1$ will be used.  For each fixed tuple
length there are finitely many such types.  Although a complete type is
an infinite collection of formulas, its class of realizations in these
two finite structures can be isolated by one finite formula.

\begin{lemma}[Extension counts]
\label{hf:extension-counts}
Let $P,P'$ be tuples of the same length $a$, in two structures chosen
from $B_0,B_1$, allowing the two structures to be the same.
Suppose that $P,P'$ have the same $C^M$ type.
If $a+d\leq M$, then for every realized type $\rho$ of length $a+d$,
the numbers of tuples $\gamma,\gamma'$ of length $d$ satisfying
\[
 \operatorname{tp}_M(P,\gamma)=\rho,
 \qquad
 \operatorname{tp}_M(P',\gamma')=\rho
\]
are equal.  Equal types also give equal types on every projected or
reordered subtuple.
\end{lemma}

\begin{proof}
Projection and reordering follow by renaming variable positions:
an injection from the selected positions to the original positions
extends to a permutation of the $M$ variable names.  This renames bound
variables as well, so it avoids capture.  Repeated atom values cause
no difficulty, because equality formulas record them.

For $d=1$, choose a tuple realizing $\rho$.  For every tuple in either
finite structure whose type differs from $\rho$, choose a formula
distinguishing it from this representative, and negate that formula
if necessary so that the representative satisfies it.  The conjunction
of these finitely many formulas, denoted $\theta_\rho$, holds exactly
on the tuples of type $\rho$ in the two structures.  It has free
variables among the first $a+1$ positions and uses the same $M$ variable
names.  If there are no other types, take a tautology.
If $P$ has exactly $h$ extensions realizing $\rho$, then
\[
       \exists^{=h}x_{a+1}\,
       \theta_\rho(x_1,\ldots,x_a,x_{a+1})
\]
holds at $P$ and hence at $P'$.  This includes $h=0$.
For empty prefixes in different structures, agreement of the types
means exactly the assumed agreement on $C^M$ sentences.

For $d>1$, let $\sigma$ be the projection of $\rho$ to the first
$a+d-1$ positions.  By induction the numbers of extensions from $P,P'$
to $\sigma$ agree.  The one-coordinate case, applied both within one
structure and between the two structures, shows that each realization
of $\sigma$ has the same number of extensions to $\rho$.
Multiplying the two finite counts proves the assertion.
\end{proof}

In particular, every type realizable over one prefix is realizable over
any prefix with the same type.  We will use this consequence to transfer
witnesses, and the full counting assertion to transfer cardinalities.

\subsection{Finite expressions for supported objects}

Call a tuple in $B_i^s$ a \emph{molecule}.  We now define a common syntax,
independent of $i$, for objects represented over a molecule.

There are atomic forms $c_1,\ldots,c_s$.  A set form is a finite set
of pairs $(\psi,\rho)$, where $\psi$ is an already constructed form and
$\rho$ is a realized $C^M$ type of length $2s$.
Atomic forms and set forms are distinct syntactic constructors.
They have finite depth: an atomic form has depth zero; an empty set form
has depth one; and a nonempty set form has depth one plus the largest
depth of any of its constituent forms.
For $\alpha\in B_i^s$, define evaluation recursively by
\begin{align}
 c_j*\alpha&=\alpha_j, \label{hf:atomic-form}\\
 \phi*\alpha&=
 \{\psi*\beta:(\psi,\rho)\in\phi,\ \beta\in B_i^s,\
             \operatorname{tp}^{B_i}_M(\beta,\alpha)=\rho\}
       \quad\text{for a set form $\phi$.}
       \label{hf:set-form}
\end{align}
This is a finite set because both $\phi$ and $B_i$ are finite.
The empty set form evaluates to $\emptyset$.

\begin{lemma}[Representation by forms]
\label{hf:representation}
Every value $\phi*\alpha$ lies in $D_i$ and is supported by $\alpha$.
Every object of $D_i$ has a representation of this kind.
Furthermore, for each $g\in K_i$,
\[
                         g(\phi*\alpha)=\phi*(g\alpha).
\]
\end{lemma}

\begin{proof}
Automorphisms preserve $C^M$ types.  Induction on the form proves the
displayed identity, using the change of witness molecule
$\beta\mapsto g\beta$ in Equation~\eqref{hf:set-form}.
It follows that $\alpha$ supports $\phi*\alpha$.
The same induction shows that every member of a set-form value, and
all of its recursive members, has a support of length $s$.
Thus every form value belongs to $D_i$.

Conversely, use induction on the hereditary construction of an object
$x\in D_i$.  If $x$ is an atom, choose a molecule with first coordinate
$x$ and use $c_1$.  If $x$ is a set, choose a supporting molecule
$\alpha$.  For every $y\in x$, induction gives
$y=\psi_y*\beta_y$.  Set
\[
 \phi=\{(\psi_y,\operatorname{tp}^{B_i}_M(\beta_y,\alpha)):y\in x\}.
\]
Every $y\in x$ belongs to $\phi*\alpha$.  Conversely, a member of
$\phi*\alpha$ has the form $\psi_y*\beta$ with
\[
 \operatorname{tp}^{B_i}_M(\beta,\alpha)
      =\operatorname{tp}^{B_i}_M(\beta_y,\alpha).
\]
Homogeneity for tuples of length $2s$ gives $g\in K_i$ sending
$(\beta_y,\alpha)$ to $(\beta,\alpha)$.
Thus $g$ fixes $\alpha$, so $gx=x$, and equivariance gives
$\psi_y*\beta=gy\in x$.  Hence $\phi*\alpha=x$.
For the empty set, $\phi$ is the empty set form.
\end{proof}

There is no depth bound in this representation.
The next lemma shows that equality and membership of form values depend
only on the joint atom type of their molecules.
Its variable requirement is independent of the depths of the forms.

\begin{lemma}[Transfer of atomic facts]
\label{hf:atomic-transfer}
Let $\phi,\psi$ be forms, and let $(\alpha,\beta)$ and
$(\alpha',\beta')$ be pairs of molecules having the same $C^M$ type,
in the same or different base structures.  Then
\begin{align*}
 \phi*\alpha=\psi*\beta
   &\ \Longleftrightarrow\
       \phi*\alpha'=\psi*\beta',\\
 \phi*\alpha\in\psi*\beta
   &\ \Longleftrightarrow\
       \phi*\alpha'\in\psi*\beta'.
\end{align*}
Atomhood and all base relations also agree on the corresponding values.
\end{lemma}

\begin{proof}
For equality, induct on the sum of the depths of $\phi,\psi$.
If both forms are atomic, equality compares two molecule coordinates
and is recorded by the joint type.  If exactly one is atomic,
equality is false on both sides.

Suppose now that both are set forms and
$\phi*\alpha=\psi*\beta$.
Take a member of $\phi*\alpha'$, witnessed by
$(\theta,\rho)\in\phi$ and $\gamma'$ such that
$\operatorname{tp}_M(\gamma',\alpha')=\rho$.
By Lemma~\ref{hf:extension-counts}, the matching prefix $(\alpha,\beta)$
can be extended by $\gamma$ so that
\[
 \operatorname{tp}_M(\alpha,\beta,\gamma)
    =\operatorname{tp}_M(\alpha',\beta',\gamma').
\]
Thus $\theta*\gamma\in\phi*\alpha=\psi*\beta$.
Choose its membership witness $(\xi,\sigma)\in\psi$ and $\zeta$ with
\[
 \operatorname{tp}_M(\zeta,\beta)=\sigma,
 \qquad
 \theta*\gamma=\xi*\zeta .
\]
Extend the primed prefix of length $3s$ once more, obtaining $\zeta'$
such that
\[
 \operatorname{tp}_M(\alpha,\beta,\gamma,\zeta)
    =\operatorname{tp}_M(\alpha',\beta',\gamma',\zeta').
\]
This uses $4s\leq M$.
Projection preserves the membership-witness types.
Since $\theta$ and $\xi$ have smaller depths than $\phi$ and $\psi$,
the induction hypothesis gives
$\theta*\gamma'=\xi*\zeta'$.
Therefore the chosen member lies in $\psi*\beta'$.
Interchanging $\phi,\psi$ proves the other inclusion; interchanging
the primed and unprimed data proves the reverse implication.
Empty set forms are included because their inclusions are vacuous.

For membership, if the right-hand form is atomic then membership is
false.  Otherwise a witness for $\phi*\alpha\in\psi*\beta$ consists of
$(\theta,\rho)\in\psi$ and a molecule $\gamma$ with
\[
 \operatorname{tp}_M(\gamma,\beta)=\rho,\qquad
 \phi*\alpha=\theta*\gamma .
\]
Extend $(\alpha',\beta')$ to match
$(\alpha,\beta,\gamma)$, using $3s$ positions.
The equality transfer already proved transfers the equality to this
witness, and hence transfers membership.  The converse is symmetric.

A form evaluates to an atom exactly when it is atomic.
Base relations are false if any argument is a set.
Otherwise, each argument is a specified molecule coordinate, so the
truth of a base relation is recorded by the joint type.
\end{proof}

\subsection{Bijections after retaining parameters}

The preceding lemmas transfer individual values.  Exact counting
requires matching all possible values simultaneously after the other
variables have been assigned.

Call two assignments to the same $r$ variable names
\emph{represented alike} if, for common forms $\phi_1,\ldots,\phi_r$,
their values have representations
\[
 x_j=\phi_j*\alpha_j,\qquad
 x'_j=\phi_j*\alpha'_j \quad(1\leq j\leq r),
\]
and the concatenations
\[
 P=(\alpha_1,\ldots,\alpha_r),\qquad
 P'=(\alpha'_1,\ldots,\alpha'_r)
\]
have the same $C^M$ type.
Dropping assigned positions preserves this property by projection.
The empty assignments are represented alike by the sentence-equivalence
hypothesis.

\begin{lemma}[Extension bijections]
\label{hf:bijection}
For represented-alike assignments with $r\leq m-1$,
there is a bijection $f:D_0\to D_1$ such that adjoining $x,f(x)$
gives represented-alike assignments for every $x\in D_0$.
Moreover, $f(x_j)=x'_j$ for every retained position $j$.
\end{lemma}

\begin{proof}
Fix the representing tuples $P,P'$.
A label is a pair $(\phi,\tau)$ consisting of a form and a realized
type of length $(r+1)s$.  Its value set on the first side is
\[
 X_{\phi,\tau}(P)=
 \{\phi*\gamma:\gamma\in B_0^s,\
                 \operatorname{tp}_M(P,\gamma)=\tau\};
\]
define $X_{\phi,\tau}(P')$ in $B_1$ similarly.
Lemma~\ref{hf:extension-counts} shows that these value sets are either
both empty or both nonempty.

Write $K_{0,P}$ for the pointwise stabilizer of the coordinates of $P$
in $K_0$, and $K_{1,P'}$ for the corresponding stabilizer in $K_1$.
Every nonempty $X_{\phi,\tau}(P)$ is a single $K_{0,P}$-orbit.
Indeed, homogeneity for
\[
                (r+1)s\leq\max\{2,m\}s
\]
shows that the molecules of the specified type over $P$ form a single
$K_{0,P}$-orbit.  Equivariance of form evaluation makes their value set
one orbit as well.  The same holds on the second side.
In particular two nonempty label sets on either side are equal or
disjoint.

We compare first their overlaps and then their cardinalities.
Suppose two labels overlap on the first side, with molecules
$\gamma,\zeta$ witnessing equal form values.
Extend $P'$ to match the type of $(P,\gamma,\zeta)$.
This uses
\[
                     (r+2)s\leq(m+1)s\leq M.
\]
Projection preserves each label condition, and
Lemma~\ref{hf:atomic-transfer} preserves equality of the two values.
Thus the labels overlap on the second side.
The converse follows by interchanging the sides.

Equal numbers of molecules do not yet give equal numbers of represented
objects: several molecules may evaluate to the same value. We will show
that evaluation has constant fibre size on each label and that this size
agrees in the two structures.

Fix a nonempty label $(\phi,\tau)$.
Let $n$ be the number of its molecules over $P$; the corresponding
number over $P'$ is also $n$, by Lemma~\ref{hf:extension-counts}.
The map
\[
                 \gamma\longmapsto\phi*\gamma
\]
from these molecules onto $X_{\phi,\tau}(P)$ has fibres of a common
positive size.  To see this, choose preimage molecules of any two
values.  An element of $K_{0,P}$ taking the first molecule to the
second maps the first fibre bijectively onto the second by equivariance.

Choose molecules $\gamma,\gamma'$ realizing $\tau$ over $P,P'$.
The fibre containing $\gamma$ has size
\[
 h=
 \bigl|\{\zeta\in B_0^s:
       \operatorname{tp}_M(P,\zeta)=\tau,\
       \phi*\gamma=\phi*\zeta\}\bigr|.
\]
Partition the possible $\zeta$ by the finitely many realized types of
$(P,\gamma,\zeta)$.
On each type, the two displayed conditions have a constant truth
value, within either structure and between the structures:
projection determines the first, and
Lemma~\ref{hf:atomic-transfer} determines the second.
By Lemma~\ref{hf:extension-counts}, each such type has equally many
extensions over $(P,\gamma)$ and $(P',\gamma')$.
Again the required length is only $(r+2)s$.
Summing these finite extension counts shows that the corresponding
fibre on the second side has the same size $h$.
Consequently
\[
                    |X_{\phi,\tau}(P)|
                     =\frac{n}{h}
                     =|X_{\phi,\tau}(P')|.
\]

By Lemma~\ref{hf:representation}, the label sets cover the respective
domains.  Partition the nonempty labels by overlap on the first side.
The same equivalence relation is obtained on the second side.
For each label class, its value set in $D_0$ and its value set in $D_1$
are finite sets of the same size.  Choose a bijection between them.
Different classes have disjoint value sets, so the union of these
bijections is a bijection $f:D_0\to D_1$.

This construction does not require either domain to be finite.
There are countably many labels, since there are finitely many base
types of each relevant length and every form is a finite expression.
For an explicit set-theoretic choice of the bijections, fix external
enumerations of the labels and of the two countable HF universes.
Use the first label in each class, and match the members of its two
finite value sets in enumeration order.
These external enumerations are used only in this proof.

If $x$ is matched with $x'$, the class's representative label
$(\phi,\tau)$ supplies molecules $\gamma,\gamma'$ with
\[
 x=\phi*\gamma,\quad x'=\phi*\gamma',\quad
 \operatorname{tp}_M(P,\gamma)
       =\operatorname{tp}_M(P',\gamma')=\tau .
\]
Thus the extended assignments are represented alike.
If $x=x_j$ for a retained position, projection onto the molecules for
$x_j$ and $x$, followed by Lemma~\ref{hf:atomic-transfer}, forces
$x'=x'_j$.  This also proves consistency when retained values repeat.
\end{proof}

\begin{proof}[Proof of Theorem~\ref{hf:transfer}]
We prove by induction on $C^m$ formulas that represented-alike
assignments to their free variables agree on truth.
For atomic formulas involving two distinct variables, project onto
their two molecules and apply Lemma~\ref{hf:atomic-transfer}.
Atomhood is determined by the form.
For a repeated variable, $x=x$ is always true and $x\in x$ always
false in $\HF(B_i)$; a diagonal base relation is false on a set and is
otherwise determined by a coordinate in the molecule type.
Unary base relations are treated in the same way.
Boolean connectives preserve agreement.

For a quantifier binding $x$, retain only the assignments to the other
free variable names, dropping the previous value of $x$ if present.
There are at most $m-1$ retained positions.
Lemma~\ref{hf:bijection} supplies a bijection of the domains for which
every extension by $x$ is represented alike.
By the induction hypothesis, this restricts to a bijection between the
sets of values satisfying the quantified body.
It preserves existence and, for every natural number $h$, the assertion
that there are exactly $h$ satisfying values.
This remains valid when the satisfying sets are infinite: a bijection
preserves finiteness and every finite cardinality.
Universal quantification follows as well, or follows by negation.
The empty assignments are represented alike, which proves the theorem.
\end{proof}

The homogeneity hypothesis was used only for representation with two
molecules and for stabilizer orbits with at most $m$ molecules.
Four molecules suffice to transfer equality, and at most $m+1$
suffice to compare label overlaps and fibres.
Neither requirement grows with hereditary rank.
The domain bijections are mathematical witnesses to counting
equivalence; no definability of them is asserted or needed.
